\let\cleardoublepage\clearpage
\chapter{文献综述}
\section{安全强化学习研究现状}

在自动驾驶、服务机器人以及多机器人协同等真实应用场景中，安全性通常是系统能否部署的前提。机器人在运行过程中需要面对动态障碍物、传感器噪声、模型误差以及环境变化等不确定因素，即使策略在仿真环境中表现良好，进入真实系统后也可能因为状态估计偏差或外部扰动而产生不安全行为。因此，如何在策略学习和执行过程中持续保持安全约束，是机器人强化学习研究中长期存在的关键问题。

强化学习（Reinforcement Learning, RL）能够通过与环境交互学习控制策略，在处理复杂决策和自适应控制任务时具有一定优势~\cite{gu2024saferl_survey}。但这类方法通常依赖试错机制来改进策略，训练早期的探索行为并不稳定，在机器人系统中可能带来碰撞、越界或任务失败等风险。即使策略训练完成后具有较高回报，也不能直接说明其在整个学习过程或真实部署阶段始终满足安全要求。

在这一背景下，安全强化学习（Safe Reinforcement Learning, SafeRL）逐渐发展。该类方法通常将问题建模为约束马尔可夫决策过程（Constrained Markov Decision Process, CMDP）~\cite{altman1999cmdp}，在最大化期望回报的同时引入安全成本或状态约束。根据约束建模方式的不同，近期综述通常将相关方法划分为基于代价的方法和基于状态约束的方法~\cite{gu2024saferl_survey, liu2024safe_survey}。

\subsection{基于代价的方法}

基于代价的方法一般通过在优化目标中加入安全代价或惩罚项来约束策略行为。PID Lagrangian~\cite{stooke2020responsive}利用 PID 控制器调节拉格朗日乘子，使策略在奖励优化和成本约束之间进行动态平衡；IPO~\cite{liu2022ipo}采用对数屏障惩罚，将约束违背转化为优化过程中的内部惩罚；软屏障方法~\cite{wang2023soft_barrier}则从概率角度维持策略的可行性。这类方法通常不依赖特定机器人结构，具有较好的通用性，但其安全约束多属于软约束形式，实际违约率容易受到惩罚系数、乘子更新方式以及训练状态的影响，因此很难在所有阶段稳定实现零代价性能~\cite{liu2024safe_survey}。

在约束累积成本框架下，算法通常将碰撞次数或其他风险指标作为成本项，并为其设置上限。策略更新时，算法并不是直接禁止所有不安全动作，而是在收益和成本之间寻找折中解。因此，这类方法更适合任务性能与约束满足之间存在可接受权衡的场景。为了实现这一目标，相关研究常采用信赖域方法、内点法或拉格朗日松弛等约束优化技术，在策略改进过程中限制安全成本的增长。例如，CPO（Constrained Policy Optimization）~\cite{achiam2017constrained} 和 PCPO~\cite{yang2020projection} 通过引入信赖域约束，使策略更新幅度保持在安全成本允许的范围内；RCPO~\cite{tessler2018reward} 则将约束项转化为奖励函数中的惩罚项，并通过交替更新策略和乘子来逼近约束目标。除此之外，机会约束方法从概率角度限制智能体进入不安全状态的可能性，使风险发生概率低于预设阈值~\cite{ono2015chance, wang2023enforcing, chen2024probabilistic}。

总体来看，基于代价的方法已经形成了较为完整的研究体系，在不同任务和算法框架中均有较多应用。这类方法的安全性通常依赖对成本函数和惩罚参数的合理设定，约束满足也往往是统计意义或期望意义上的结果。当任务要求机器人在训练和执行过程中尽量避免任何碰撞或越界时，仅依靠代价惩罚并不足以提供严格保证。尤其在多机器人编队导航任务中，单个智能体的局部违约可能引发队形破坏或连锁碰撞，因此训练过程中的持续安全性仍然是这类方法需要进一步解决的问题。

\subsection{基于状态约束的方法}

与基于代价的方法不同，基于状态约束的方法更关注系统状态是否始终保持在安全集合内。该类方法通常通过限制动作空间或修正策略输出来避免系统进入危险状态。控制屏障函数（Control Barrier Functions, CBFs）~\cite{dawson2023safe_cbf, as2022cbfrl} 是其中较有代表性的一类方法，其核心思想是在控制输入层面对动作进行修正，使系统状态保持在前向不变的安全集内。CBF 能够提供较强的安全解释，但在复杂环境中构造合适的屏障函数并不容易，尤其当障碍物分布、动力学模型或观测条件发生变化时，屏障函数的设计和验证都会变得更加困难。Recovery RL~\cite{thananjeyan2021recovery} 则通过离线方式学习恢复区域，当智能体接近不安全边界时切换到恢复策略，从而降低进入危险状态的风险，但这种方法的安全效果很大程度上取决于恢复边界学习是否准确。

另一类研究通过约束状态可行性来间接限制策略行为。由于动作会引起状态转移，当策略选择的动作能够使系统始终停留在可行区域内，策略就可以被认为是安全的。部分研究采用时序逻辑对探索行为进行验证，以保证动作序列满足预设的安全规范~\cite{li2019temporal}；部分方法利用 Lyapunov 函数构造能量边界，将智能体动作限制在稳定区域内，避免系统演化到危险状态~\cite{chow2018lyapunov, chow2019lyapunov, koller2018learning}。这类方法通常需要对状态空间、动力学模型或约束形式作出明确假设，面对未知环境或复杂传感噪声时，泛化能力会受到限制。

高斯过程也被用于探索阶段的安全状态空间估计。相关方法通过建模环境不确定性，在策略探索过程中逐步扩展已验证的安全区域~\cite{turchetta2016safe, berkenkamp2015safe, sui2015safe, wachi2018safe}。例如，SafeMDP~\cite{turchetta2016safe} 使用高斯过程对环境进行建模，并借助置信区间评估状态安全性，使探索过程被限制在已经验证的安全状态范围内。\cite{berkenkamp2015safe} 则将高斯过程模型的在线更新嵌入鲁棒控制框架，通过数据驱动方式逐步减少系统不确定性，在保持稳定性的同时提升控制性能。此类方法能够较好地刻画不确定性，但在高维状态空间或多机器人场景中，计算开销和样本需求往往较高。

近年来，也有研究尝试使用约束流形构建安全空间，通过调整机器人动作使其沿约束允许的方向运动~\cite{liu2022robot}。这类方法将安全约束转化为几何结构，为动作修正提供了更直接的表达方式，但现有方法大多仍面向已知且相对固定的环境，面对动态障碍物、感知噪声和部分可观测条件时，其适用范围仍然有限。ATACOM~\cite{liu2022atacom, liu2023safe_manipulation} 进一步提出利用零空间投影将动作映射到约束流形的切空间，从而为机器人操作任务提供可证明的约束满足性。该方法为安全动作修正提供了清晰的几何解释，但它通常假设系统在每一时刻都能够获得精确状态信息和已知障碍物位置；在传感器噪声明显、环境信息不完整或多机器人协同约束不断变化的场景中，这一假设往往难以满足。因此，如何在保留约束流形安全修正优势的同时，引入状态估计、可达性分析和多机器人几何约束机制，仍然具有进一步研究的必要。

\section{多智能体强化学习研究现状}
在多智能体中，机器人不仅需要满足自身的安全约束并最大化回报，还需确保团队行为的整体安全性，平衡安全约束与奖励优化 \cite{chen2024duality}。
现有多智能体安全强化学习方法中，MACPO \cite{gu2023safe} 和 MAPPO-Lagrangian \cite{gu2023safe} 具有一定代表性。MACPO 在策略更新过程中引入信赖域思想，通过限制个体策略的更新幅度来降低安全成本突增的风险，并结合全局约束机制维持多智能体系统的整体安全性。
MAPPO-Lagrangian 则沿用 PPO 的裁剪目标进行策略优化，同时利用拉格朗日乘子对全局约束权重进行动态调节，使策略在任务回报和安全代价之间不断调整。此类方法为多智能体安全强化学习提供了较为直接的优化框架，也便于与现有策略梯度算法结合。
这类方法通常建立在较强的优化假设之上，尤其依赖原始问题与对偶问题之间具有较好的对应关系。
在多智能体系统中，不同智能体策略相互影响，个体策略与联合行为占用度量之间往往呈现非凸关系，强对偶性难以始终成立。
当存在非零对偶间隙时，对偶问题就不能完全刻画原始约束优化问题，拉格朗日乘子得到的解也可能与原始问题的最优解不一致，甚至出现约束满足和任务性能同时难以稳定收敛的情况。
对于多机器人编队导航这类约束耦合较强的任务，仅依赖全局代价或乘子调节往往不足以精细描述机器人间距、队形保持和避障安全等几何关系，
因此仍有必要在策略学习过程中引入更直接的安全约束表达与动作修正机制。
% 本文在此流形投影范式基础上进行了三方面扩展：（i）通过离线可达性分析利用先验地图预计算环境相关安全裕度，替代实时障碍物检测与跟踪；（ii）引入数据驱动 EKF 提升噪声条件下的状态估计鲁棒性；（iii）通过奖励校准机制引导策略学习安全行为，无需外部惩罚调参。

\subsection{多智能体编队研究现状}

近年来，无人系统在巡检、救援、运输和协同作业等场景中的应用不断扩展，多智能体编队控制也因此受到较多关注。编队控制通常要求多个智能体在运动过程中形成并维持特定空间结构，同时完成目标到达、避障和协同行动等任务。已有综述将编队控制方法概括为领航-跟随（leader-follower）\cite{sui2020formation}、虚拟结构（virtual structures）\cite{yan2021formation} 和基于行为的方法（behavior-based approaches） \cite{liu2018survey}。

领航-跟随方法结构简单，便于实现，但对领航者依赖较强，一旦领航者失效，编队任务容易受到影响。虚拟结构方法将编队视为整体结构进行控制，能够保持较稳定的队形，但在队形调整和复杂环境适应方面灵活性不足。相比之下，基于行为的方法 \cite{xu2014behavior, quan2023robust} 更强调机器人根据局部观测动态调整动作，适合处理目标到达、队形保持和避障等多重要求。随着多智能体强化学习的发展，MARL 已逐渐用于编队控制任务，并取得了一定成果 \cite{sui2020formation,yan2022relative,khan2020graph}。

现有强化学习编队方法中，PPO 及其多智能体扩展应用较为广泛。文献 \cite{sadhukhan2022proximal} 将 PPO 与课程学习结合，通过机器人之间的夹角信息维持编队形状；文献 \cite{yan2022relative} 采用 MAPPO，并结合课程学习、参考目标和相对位置信息实现编队控制。针对四旋翼编队，文献 \cite{batra2022decentralized} 基于 PPO 引入注意力机制，对相邻四旋翼状态进行编码，以减少碰撞并提升飞行性能，其思想与文献 \cite{hoshen2017vain, chen2019crowd} 中的注意力建模方法相关。文献 \cite{wang2019pattern} 则在集中训练、分散执行（Centralized Training with Decentralized Execution, CTDE）框架下，利用 PPO 实现移动机器人的静态编队控制。

除 PPO 类方法外，其他强化学习框架也被用于编队控制。文献 \cite{sui2020formation} 将模仿学习与强化学习结合，在领航-跟随框架下构建分布式编队控制器。文献 \cite{khan2020graph} 将策略梯度方法与图卷积网络（Graph Convolutional Network, GCN）结合，用于无人机静态编队控制。文献 \cite{xie2021reinforcement} 和文献 \cite{zhou2019learn} 分别基于深度确定性策略梯度（Deep Deterministic Policy Gradient, DDPG）和深度 Q 网络（Deep Q-Network, DQN），实现了海上无人艇的领航--跟随编队控制。

总体来看，现有研究已证明强化学习方法在多智能体编队控制中的可行性，但多数工作更关注队形形成和任务完成效果，对安全距离、障碍物规避以及训练过程中的约束违背考虑仍不充分。在复杂环境下，仅依赖奖励函数引导编队行为，难以保证策略输出始终满足安全要求。因此，将几何约束、安全过滤和动作修正机制引入多智能体策略学习，仍具有进一步研究价值。

\subsection{多智能体安全强化学习}
在编队控制下，安全性意味着在训练与执行过程中避免任何形式的碰撞。

MARL 的核心目标是在多种交互场景（包括协作、竞争及混合环境）中协调多个智能体之间的行为。
安全多智能体强化学习（Safe MARL）近年来已成为研究热点之一 \cite{garg2024learning}。

文献 \cite{cai2021safe} 将 CBF 与多智能体深度确定性策略梯度（MADDPG）\cite{lowe2017multi} 相结合，实现了两智能体在导航过程中无碰撞运动，其中 CBF 通过构造满足安全约束的控制律来保证系统安全。文献 \cite{zhang2023spatial} 研究了注意力模块与基于 CBF 的分布式安全屏蔽机制在多自动驾驶车辆场景中降低碰撞率的有效性。

不同于上述方法，文献 \cite{sheebaelhamd2021safe} 在 MADDPG 框架中引入训练得到的集中式神经网络作为安全层，并在协作导航任务中显著降低了碰撞发生率。文献 \cite{elsayed2021safe} 则基于确定性有限自动机（Deterministic Finite Automaton, DFA）提出了集中式与分布式相结合的安全屏蔽策略，以确保协作导航任务中的智能体安全，其中 DFA 通过定义有限状态集合来刻画安全约束。

文献 \cite{zhang2019mamps} 提出了一种预测式安全屏蔽方法，并与 MADDPG 相结合以保证多智能体协作导航的安全性，并在简化的多智能体粒子环境 \cite{mordatch2018emergence} 中验证了其有效性。

基于行为的编队控制问题中，由于机器人需要在空间上保持紧密相对位置，碰撞风险比传统协作导航场景更高。已有研究多集中于协作导航任务，该类任务通常可通过训练单一策略并部署至多个机器人来实现 \cite{everett2021collision, brito2021go, vinod2022safe}。然而，行为驱动的编队控制需要机器人在训练过程中持续交互，因此更适合采用多智能体强化学习方法。文献 \cite{zhang2022barrier} 利用基于强化学习的模型预测控制方法，在保证安全的前提下实现了移动机器人的领航–跟随编队控制。该方法通过为每个机器人设定目标位置来维持编队结构。

\section{Hamilton–Jacobi 可达性分析}

Hamilton–Jacobi（HJ）可达性分析~\cite{bansal2017hj_reachability}通过求解 HJ 变分不等式（HJ-VI）计算反向可达管（Backward Reachable Tubes, BRTs），从而基于受控不变集提供形式化的安全保证。近期研究致力于提升其可扩展性：例如基于神经网络的近似方法~\cite{hsu2021safety, so2023rtreach}能够在更高维空间中进行可达性计算。然而，大多数方法仍需在线计算或依赖环境动力学模型。本文则采用离线、数据驱动的方式：将 HJ-VI 离散化为可行性 Bellman 算子，并利用期望分位回归~\cite{kostrikov2022iql}从采集轨迹数据中求解，从而在无需显式动力学模型的情况下近似 BRT 值函数。结合静态障碍物先验地图，该方法实现了安全裕度计算与部署阶段的解耦，适用于资源受限的机器人平台。

% \subsection{多任务几何约束融合}
% 为解决多任务几何约束融合问题，RMPflow 提出了一种基于黎曼几何的运动策略组合框架。RMPflow 将每个任务空间中的控制策略表示为 Riemannian Motion Policy (RMP)，并通过雅可比映射在不同空间之间进行一致变换与组合，从而实现多个任务策略的解析融合~\cite{cheng2018rmp}。

% RMPflow 已被应用于机械臂操作、移动机器人导航以及人机协作等场景，在复杂约束环境中展现出良好的稳定性与鲁棒性~\cite{cheng2018rmp, li2021rmp2}。然而，该方法通常依赖于精确系统模型与人工设计的任务策略，其在高度不确定环境中的适应能力有限。此外，RMPflow 本身属于模型驱动方法，缺乏从交互数据中学习高层决策策略的能力，这在多机器人协同探索与动态环境任务中成为主要限制。

% 近年来，研究者开始探索将强化学习与几何约束控制相结合的框架，以同时获得学习方法的适应性与几何控制的稳定性。例如，通过在强化学习策略外层叠加几何安全控制器，或采用分层控制结构，使高层策略负责任务决策，而低层控制器保证运动学与动力学约束的满足~\cite{deisenroth2013gaussian,}. 这种融合方法被认为是实现复杂环境中安全多机器人协作的重要方向。

% 本文提出将 RMPflow 与强化学习进行分层融合，在保证几何安全约束的同时提升多机器人协作策略的环境适应能力。


% 黎曼运动策略（Riemannian Motion Policies，简称 RMPs）是一种定义在流形上的策略的数学表示方法，而 RMPflow 是一种用于组合多个 RMP 的递归算法。

% 一个机器人或一组机器人，其配置空间 C 是一个光滑的 d 维流形。假设 C 存在一个全局的广义坐标 $q : C \to \mathbb{R}^d$ ，系统可以通过反馈线性化，使其可以直接通过广义加速度进行控制，即 $\ddot{q} = u(q, \dot{q})$。我们称 u 为策略或控制器，$(q, \dot{q})$为系统状态。

% RMPflow  将一个任务拆解为一系列子任务，例如避开障碍物、到达目标、跟踪轨迹等，即任务空间 T 由多个子任务空间组成的集合，每个空间对应一个子任务。假设任务空间 T 与配置空间 C 可以通过一个平滑的任务映射 $\psi : C \to T$ 相互关联。RMPflow 的目标是在配置空间 C 中生成控制策略 u，使得机器人在任务空间 T 中做出期望的行为。

% \section{黎曼运动策略（RMPs）}
% 黎曼运动策略（Riemannian Motion Policies, RMPs） 是一种定义在流形上的控制策略表示方法。设 $\mathcal{M}$ 是一个维度为 m 的流形，其广义坐标为 $x \in \mathbb{R}^m$。RMP 在流形 $\mathcal{M}$ 上可以通过标准形式和自然形式两种形式进行表示。

% RMP 的标准形式表示为一对 $(a, M)^{\mathcal{M}}$，其中：$a : (x, \dot{x}) \mapsto a(x, \dot{x}) \in \mathbb{R}^m$ 是期望加速度（即控制输入）；$M : (x, \dot{x}) \mapsto M(x, \dot{x}) \in \mathbb{R}^{m \times m}_+$ 是惯性矩阵（或权重矩阵），用于在与其他 RMP 合并时表示该策略的重要性。

% 根据标准形式，我们可以得到 RMP 的自然形式，表示为一对 $[f, M]^{\mathcal{M}}$，其中 $f = M a$ 是期望力。引入自然形式主要是出于计算上的便利。

% 在流形 $\mathcal{M}$ 上的 RMP 通常来自一类称为几何动力系统（Geometric Dynamical Systems，GDSs） 的系统。GDS 是经典的简单机械系统（Simple Mechanical Systems, SMSs）的一种推广。在 GDS 中，动能度量 G 是配置和速度的函数，即 $G(x, \dot{x}) \in \mathbb{R}^{m \times m}_+$。这意味着动能可以依赖于运动方向，这在避障等应用中是有用的。GDS 的动力学形式如下：
% \begin{equation}
%     (G(x, \dot{x}) + \Xi_G(x, \dot{x})) \ddot{x} + \xi_G(x, \dot{x}) = -\nabla_x \Phi(x) - B(x, \dot{x}) \dot{x} 
% \end{equation}
% 其中：$B : \mathbb{R}^m \times \mathbb{R}^m \to \mathbb{R}^{m \times m}_+$ 是阻尼矩阵；$\Phi : \mathbb{R}^m \to \mathbb{R}$ 是势函数；曲率项 $\Xi_G$ 和 $\xi_G$ 由度量 G 所诱导，具体表达为：
% \begin{equation}
% \begin{aligned}
%     \Xi_G(x,\dot{x}) &:= \frac{1}{2} \sum_{i=1}^{m} \dot{x}_i, \partial_{\dot{x}} g_i(x, \dot{x}) \\
%     xi_G(x, \dot{x}) &:= \partial_x G(x, \dot{x}) \dot{x} - \frac{1}{2} \nabla_x ( \dot{x}^\top G(x, \dot{x}) \dot{x} )
% \end{aligned}
% \end{equation}
% 其中：$g_i$ 表示 G 的第 i 列；$x_i$ 表示 x 的第 i 个分量；$\partial_x G(x, \dot{x}) \dot{x} = [\partial_x g_i(x, \dot{x}) \dot{x}]_{i=1}^m$，即将各列拼接为矩阵。

% 给定上述 GDS 动力学，我们可以自然地构造出对应的 RMP：其标准形式为：$a = \ddot{x}, \quad M(x, \dot{x}) = G(x, \dot{x}) + \Xi_G(x, \dot{x})$因此，速度相关的度量 G 在与其他 RMP 合并时，提供了一个速度相关的重要性权重 M。

% \section{ RMPflow}

% RMPflow  是一种用于在构型空间上生成控制策略的算法，它依赖于各个子任务的 RMP（例如避障、到达目标点等）。给定机器人在构型空间中的状态信息，以及针对各个子任务独立设计的控制器（即 RMPs），RMPflow 通过组合这些控制器，在构型空间上生成控制输入。RMPflow 引入了以下两个关键结构： 
% \begin{itemize}
%     \item 一种数据结构：RMP 树（RMP-tree），用于描述任务映射 ψ 的结构； 
%     \item 一组算子：RMP 代数（RMP-algebra），用于在 RMP 树中传播信息。
% \end{itemize}

% RMP 树是一个有向树。树中每个节点 u 都关联一个定义在流形 $\mathcal{M}$ 上的状态 $(x, \dot{x})$，以及一个 RMP $(f_u, M_u)^{\mathcal{M}}$。每条边 e 则带有一个从父节点到子节点的光滑映射 $\Psi_e$。例如，如图~\ref{fig:rmp树}所示的是一个典型的 RMP 树。树的根节点 r 对应机器人在构型空间中的状态$(q, \dot{q})$，以及该空间上的控制策略 $(f_r, M_r)^C$。每个叶节点$l_k$表示一个子任务，其控制策略为定义在子任务空间 $T_k$上的 RMP $(f_{l_k}, M_{l_k})^{T_k}$。
% \begin{figure}
%     \centering
%     \includegraphics[width=1\linewidth]{硕士毕业论文/毕业论文-safeMARL/images/rmp.png}
%     \caption{RMP树示例}
%     \label{fig:rmp树}
% \end{figure}
% RMP代数 包括三个基本操作符：前向传播（pushforward）、反向传播（pullback） 和解析（resolve）。具体作用过程如下：
% 设某节点 u 有 N 个子节点，分别记为 ${v_j}_{j=1}^N$，边记为 ${e_j}_{j=1}^N$，如图~\ref{fig:rmp树}所示。假设 u 所在流形为 M，每个子节点 $v_j $所在流形为 $N_j$，则各个操作符作用如下：
% 前向传播将状态从父节点 u 传播到其子节点 ${v_j}_{j=1}^N$。 给定 u 的状态 $(x, \dot{x})$，对于每个子节点 $v_j$，通过映射 $\psi_{e_j}$ 计算其状态为：
% \begin{equation}
%     (y_j, \dot{y}_j) = (\psi_{e_j}(x), J_{e_j}(x) \dot{x})
% \end{equation}
% 其中 $J_{e_j} = \frac{\partial \psi_{e_j}}{\partial x}$ 是映射 $\psi_{e_j}$ 的雅可比矩阵。

% 反向传播将子节点${v_j}_{j=1}^N$的 RMP 聚合为父节点 u 的 RMP。给定每个子节点的 RMP $(f_{v_j}, M_{v_j})^{N_j}$，父节点的 RMP $(f_u, M_u)^M$ 计算如下：
% \begin{equation}
%     f_u = \sum_{j=1}^{N} J_{e_j}^T (f_{v_j} - M_{v_j} \dot{J}_{e_j} \dot{x}), \quad M_u = \sum_{j=1}^{N} J_{e_j}^T M_{v_j} J_{e_j}
% \end{equation}

% 解析操作将自然形式（natural form）的 RMP 转换为标准形式（canonical form）。 给定自然形式 RMP $(f_u, M_u)^M$，标准形式为：$(a_u, M_u)^M, \quad a_u = M_u^{\dagger} f_u$其中 $\dagger$ 表示 Moore-Penrose 伪逆。

% 在 RMP 树结构确定之后，RMPflow 会按照前向传播、叶节点评估、反向传播以及解析输出的计算流程生成最终控制策略，具体过程如下：

% 在前向传播阶段，系统以根节点为起点，将当前机器人的状态沿任务树逐层向下传递，该过程通过递归方式完成。pushforward 操作将配置空间中的位姿与速度映射到各个任务空间，使不同层级的节点都能够获得与自身相关的状态信息。由于叶节点通常定义在不同的空间中，或者描述距离关系，或者描述相对位姿，或者直接与目标位置相关，因此需要借助雅可比矩阵进行变量变换，将全局状态转换为各任务所需的局部变量及其变化率，使各叶节点能够在各自坐标系下进行计算。

% 在叶节点评估阶段，系统根据节点所对应的任务目标构造自然形式的 RMP $(f_{l_k}, M_{l_k})$。该形式一般由几何动力系统给出，在任务空间内通过势函数与阻尼项刻画系统的演化趋势，从而生成具有物理意义的广义力 $f_{l_k}$ 与度量矩阵 $M_{l_k}$。广义力反映任务期望的加速度方向，而度量矩阵描述任务在当前状态下的作用强度及方向权重。当机器人接近约束边界时，与安全相关的任务度量会增大，使其在后续融合中占据更高权重；而用于引导系统接近目标的任务，其度量变化相对平缓。

% 在反向传播过程中，系统从所有叶节点出发，通过 pullback 操作将自然形式的 RMP 逐层上传至父节点，并在传递过程中完成合并。不同任务空间中得到的广义力与度量矩阵被统一映射回根节点所在的配置空间，并依据各自的度量进行加权叠加。该过程保证各任务在几何上的一致性，使其能够在同一坐标框架下融合。当某一任务的度量值较大时，其在合成结果中的影响随之增强，而其他任务的作用相对减弱。

% 当所有叶节点的信息回传至根节点后，系统对汇总得到的自然形式 $(f_r, M_r)$ 进行解析，将其转换为可直接用于控制的标准形式。具体做法是求解加权最小二乘意义下的最优加速度：
% \begin{equation}
% a_r = M_r^{\dagger} f_r
% \end{equation}

% 该求解过程等价于在统一的黎曼度量下寻找整体上最符合各项任务约束的加速度解。与基于硬优先级的传统控制方式相比，该方法依赖度量矩阵连续调节各任务的影响程度，使权重变化保持平滑，从而减小任务切换带来的突变与振荡。

% 在执行阶段，机器人按照 $\ddot q = u = a_r$ 的形式施加控制输入。对于移动机器人系统，配置空间中的加速度还需进一步映射为底层驱动可执行的速度或角速度指令，经此转换后控制量才能作用于执行机构。通过上述流程，系统能够在多任务与多约束条件下持续生成平滑且稳定的控制指令，并在满足物理约束的前提下协调各类控制目标。

\section{本章小结}

本章系统梳理了安全强化学习、多智能体强化学习与多机器人编队控制的研究进展。现有安全强化学习方法主要分为基于代价约束与基于状态约束两类：
前者具有通用性但通常难以实现严格零违约，后者能够提供更强的安全保证，
但常依赖精确状态估计、在线环境感知及较强模型假设。在多智能体场景中，现有方法虽在约束优化方面取得进展，
但仍面临强对偶性难以满足、策略协同复杂度高以及训练过程安全性不足等问题。

针对上述不足，本章进一步分析了编队任务中“结构保持 + 避障安全 + 目标到达”的耦合挑战，
并综述了 Hamilton--Jacobi 可达性分析在形式化安全保证方面的优势及其计算代价问题。
基于此，本文后续章节将以安全可达空间为先验约束基础，结合约束流形动作修正与数据驱动状态估计，
构建兼顾安全性、可部署性与协同性能的多智能体安全强化学习方法。